\begin{figure}[t]
\centering
\begingroup
\definecolor{llmred}{RGB}{249,220,220}%
\definecolor{llmredline}{RGB}{201,74,74}%
\definecolor{gpblue}{RGB}{214,231,247}%
\definecolor{gpblueline}{RGB}{58,110,165}%
\definecolor{graybox}{RGB}{238,238,238}%
\definecolor{grayline}{RGB}{120,120,120}%
\definecolor{fbgray}{RGB}{90,90,90}%
\resizebox{\textwidth}{!}{%
\begin{tikzpicture}[
    >={Latex[length=2.4mm]},
    font=\normalsize,
    llm/.style={draw=llmredline,line width=1pt,fill=llmred,rounded corners=5pt,
                align=center,inner sep=7pt,minimum height=20mm,text width=50mm},
    gp/.style={draw=gpblueline,line width=1pt,fill=gpblue,rounded corners=5pt,
                align=center,inner sep=7pt,minimum height=20mm,text width=50mm},
    proc/.style={draw=grayline,line width=0.9pt,fill=graybox,rounded corners=4pt,
                align=center,inner sep=6pt},
    lbl/.style={font=\small,align=center},
    fb/.style={font=\small\itshape,align=center,text=fbgray},
    flow/.style={->,line width=0.9pt},
    fbflow/.style={->,line width=1.0pt,fbgray},
]
 
% Top boxes
\node[llm] (llm) at (-2.9,3.7)
  {\textbf{LLM reservoir}\\[3pt]$p_{\theta,\alpha}(\cdot\mid\C_t)$};
\node[gp]  (gp)  at ( 2.9,3.7)
  {\textbf{GP surrogate: reward model}\\[3pt]$\mu_t,\sigma_t$, acquisition $\acq_t$};
 
% Tilted-search box
\node[proc,text width=78mm] (pi) at (0,0.65)
  {\textbf{acquisition-tilted search}\\[3pt]
   $\pi_t(x)\;\propto\;p_{\theta,\alpha}(x\mid\C_t)\,\exp\!\big\{\eta\,\acq_t(x)\big\}$};
 
% Candidate batch
\node[proc,text width=70mm] (cand) at (0,-1.65)
  {\textbf{draw candidate batch }$B_t=\{x_{t,1},\dots,x_{t,b}\}$};
 
% Evaluate batch
\node[proc,text width=68mm] (eval) at (0,-3.95)
  {\textbf{evaluate }$R(x)$\textbf{ for all }$x\in B_t$};
 
% Arrows: top boxes into pi
\draw[flow] (llm.south) -- (pi.north -| llm.south)
   node[lbl,midway,left=3pt] {base measure};
\draw[flow] (gp.south) -- (pi.north -| gp.south)
   node[lbl,midway,right=3pt] {acquisition function};
 
% pi -> cand -> eval
\draw[flow] (pi) -- (cand)
   node[lbl,midway,right=1pt] {inference-time search};
\draw[flow] (cand) -- (eval)
   node[lbl,midway,right=1pt] {query};
 
% Feedback loops (both labelled "feedback")
% eval -> llm  (update history)
\draw[fbflow] (eval.west) .. controls +(-3.6,0) and +(-1.8,0) .. (llm.west)
   node[fb,pos=0.28,fill=white,inner sep=2pt] {feedback}
   node[lbl,pos=0.62,fill=white,inner sep=2pt] {update history $\C_t$};
% eval -> gp  (update data)
\draw[fbflow] (eval.east) .. controls +(3.6,0) and +(1.8,0) .. (gp.east)
   node[fb,pos=0.28,fill=white,inner sep=2pt] {feedback}
   node[lbl,pos=0.62,fill=white,inner sep=2pt] {update data $\D_t$};
 
\end{tikzpicture}%
}
\endgroup
\caption{The Large Discovery Model loop. The LLM supplies the base measure
$p_{\theta,\alpha}(\cdot\mid\C_t)$ and the GP surrogate supplies the acquisition function
$\acq_t$; together they define the acquisition-tilted search distribution $\pi_t$. At each
round, inference-time search draws a candidate batch $B_t=\{x_{t,1},\dots,x_{t,b}\}$,
evaluates $R(x)$ for every $x\in B_t$, and feeds the observations back to update data $\D_t$
and history $\C_t$.}
\label{fig:loop}
\end{figure}